\documentclass[11pt,a4paper]{article}
\usepackage[margin=2.2cm,headheight=14pt]{geometry}
\usepackage[T1]{fontenc}
\usepackage[utf8]{inputenc}
\usepackage{babel}
\babelprovide[import=id, main]{indonesian}
\usepackage{lmodern}
\usepackage{microtype}
\DisableLigatures{family=tt*}
\usepackage{graphicx}
\usepackage{xcolor}
\usepackage{booktabs}
\usepackage{array}
\usepackage{tabularx}
\usepackage{float}
\usepackage{caption}
\usepackage{enumitem}
\usepackage{fancyhdr}
\usepackage{titlesec}
\usepackage{hyperref}
\usepackage{url}
\usepackage{listings}
\usepackage{seqsplit}
\usepackage{soul}
\sethlcolor{softgray}
\definecolor{blueaccent}{HTML}{2783DE}
\definecolor{softblue}{HTML}{E5F2FC}
\definecolor{darktext}{HTML}{2C2C2B}
\definecolor{softgray}{HTML}{F3F4F6}
\hypersetup{colorlinks=true,linkcolor=blueaccent,urlcolor=blueaccent,pdftitle={Tutorial GPU: Pemrograman CUDA C++},pdfauthor={Aufa Akmal Bunaya}}
\pagestyle{fancy}
\fancyhf{}
\lhead{Tutorial GPU: CUDA C++}
\rhead{Aufa Akmal Bunaya}
\cfoot{\thepage}
\setlength{\parindent}{0pt}
\setlength{\parskip}{6pt}
\titleformat{\section}{\Large\bfseries\color{darktext}}{}{0pt}{}
\titleformat{\subsection}{\large\bfseries\color{darktext}}{}{0pt}{}
\newcommand{\cmd}[1]{\hl{\texttt{#1}}}
\newcommand{\plotfig}[3]{%
  \begin{figure}[H]\centering
  \includegraphics[width=#1]{#2}
  \caption{#3}
  \end{figure}}
\lstset{basicstyle=\ttfamily\small,breaklines=true,frame=single,backgroundcolor=\color{softgray}}
\begin{document}

\begin{titlepage}
\centering
\vspace*{1.3cm}
{\Large\bfseries UNIVERSITAS GADJAH MADA\par}
\vspace{0.4cm}
{\large Departemen Ilmu Komputer dan Elektronika\par}
\vspace{2.0cm}
{\Huge\bfseries Tutorial GPU: Pemrograman CUDA C++\par}
\vspace{0.5cm}
{\large Delapan Exercise dan Mini-Project Stencil 2D: Inspeksi Device, Kernel, Vector Add, Timing, Grid-Stride, Reduksi, Matmul Tiled, dan Compute Sanitizer\par}
\vfill
\begin{tabular}{>{\bfseries}rl}
Nama: & Aufa Akmal Bunaya\\
NIM: & 23/515767/PA/22027\\
Repository GitHub: & \href{https://github.com/aufaakmalbunaya/cuda-tutorial}{github.com/aufaakmalbunaya/cuda-tutorial}
\end{tabular}
\vfill
{\large 4 Oktober 2026\par}
\end{titlepage}

\tableofcontents
\newpage

\section{Pendahuluan}
Tugas ini mengerjakan tutorial pemrograman GPU dengan CUDA C++: delapan exercise berurutan yang mencakup inspeksi device, peluncuran kernel pertama, penambahan vektor dengan manajemen memori device eksplisit, pemisahan waktu kernel vs waktu transfer, kernel grid-stride dengan tuning konfigurasi launch dan bandwidth efektif, reduksi memakai shared memory, perbandingan perkalian matriks naive vs tiled, serta diagnosis error dengan Compute Sanitizer, ditutup dengan mini-project stencil 2D lima titik. Seluruh program dikompilasi dan dijalankan secara nyata pada GPU NVIDIA; semua angka pada laporan ini adalah hasil pengukuran nyata (bukan estimasi), diringkas sebagai median dari beberapa run.

\subsection{Lingkungan dan persiapan}
VM lokal yang tersedia tidak memiliki GPU (\cmd{nvidia-smi} tidak ada, \cmd{nvcc} tidak ada, tidak ada device \cmd{/dev/nvidia*}), sehingga seluruh pekerjaan dijalankan pada notebook Kaggle dengan akselerator GPU. Spesifikasi lingkungan pada Tabel~\ref{tab:env} dan toolchain pada Tabel~\ref{tab:sw}.

\begin{table}[H]\centering
\caption{Spesifikasi GPU (hasil \cmd{nvidia-smi} dan program inspeksi).}
\label{tab:env}
\begin{tabular}{ll}
\toprule
Komponen & Nilai\\
\midrule
GPU & Tesla T4, 15360 MiB (15 GiB)\\
Compute capability & 7.5\\
Streaming multiprocessors & 40\\
Warp size & 32 thread\\
Maksimum thread per block & 1024\\
Global memory & 14.56 GiB (terpakai)\\%
Shared memory per block & 49152 byte (48 KiB)\\
Driver & 580.178.04\\
\bottomrule
\end{tabular}
\end{table}

\begin{table}[H]\centering
\caption{Toolchain kompilasi.}
\label{tab:sw}
\begin{tabular}{ll}
\toprule
Komponen & Nilai\\
\midrule
CUDA toolkit & 12.8 (V12.8.93)\\
Flag kompilasi & \cmd{-std=c++17 -O3 -lineinfo -arch=native}\\
Helper makro & \cmd{CUDA\_CHECK} pada \texttt{cuda\_check.cuh} (cek tiap API call)\\
Platform eksekusi & Notebook Kaggle, akselerator GPU\\
\bottomrule
\end{tabular}
\end{table}

\section{Exercise 1: Inspeksi GPU yang tersedia}
Berkas: \texttt{src/device\_info.cu}. Program memanggil \cmd{cudaGetDeviceProperties} dan mencetak properti device 0. Hasilnya sesuai Tabel~\ref{tab:env}: Tesla T4 dengan compute capability 7.5, 40 SM, dan memori global 14.56 GiB. Flag \cmd{-arch=native} pada nvcc 12.8 berhasil dipakai langsung pada T4 tanpa perlu menentukan \cmd{sm\_75} secara manual.

\section{Exercise 2: Kernel pertama}
Berkas: \texttt{src/hello\_gpu.cu}. Kernel \cmd{hello} menghitung indeks global \cmd{i = blockIdx.x * blockDim.x + threadIdx.x} dan mencetak \cmd{block}, \cmd{local} (threadIdx.x), dan \cmd{global} hanya bila \cmd{i < n}. Dengan \cmd{n=20} dan 8 thread per block, dibutuhkan 3 block sehingga 24 thread diluncurkan dan 20 di antaranya aktif; guard \cmd{i < n} mencegah 4 thread sisa menulis/mencetak di luar rentang.

Satu detail menarik dari output nyata: baris \cmd{block=2} tercetak \emph{sebelum} \cmd{block=0} --- urutan \cmd{printf} antar block tidak deterministik karena block dieksekusi sejajar oleh SM yang berbeda. Host mencetak konfirmasi setelah \cmd{cudaDeviceSynchronize}.

Varian modifikasi (\texttt{src/hello\_gpu\_17.cu}, \cmd{n=17}, 16 thread per block): tepat 2 block diluncurkan (32 thread), 17 aktif, sesuai prediksi \cmd{ceil(17/16) = 2}.

\section{Exercise 3: Penjumlahan vektor dengan memori device eksplisit}
Berkas: \texttt{src/vector\_add.cu}. Alur program mengikuti tutorial: \cmd{cudaMalloc} untuk tiga vektor, \cmd{cudaMemcpy} host-ke-device, launch kernel \cmd{vec\_add} dengan 256 thread per block, \cmd{cudaMemcpy} device-ke-host, verifikasi terhadap hasil CPU, lalu \cmd{cudaFree}. Untuk \cmd{n=1000003} dipakai 3907 block, \cmd{errors=0}, dengan \cmd{first=0.00} dan \cmd{last=9.25} sesuai inisialisasi \cmd{a[i]=i*0.25f}, \cmd{b[i]=i*0.5f} (yakni \cmd{c[i]=0.75*i}).

Uji ukuran tepi (rekompilasi per ukuran, sesuai permintaan tutorial) semuanya lolos:
\begin{center}
\begin{tabular}{rrrr}
\toprule
\cmd{n} & block & thread & errors\\
\midrule
1 & 1 & 256 & 0\\
255 & 1 & 256 & 0\\
256 & 1 & 256 & 0\\
257 & 2 & 256 & 0\\
1000003 & 3907 & 256 & 0\\
\bottomrule
\end{tabular}
\end{center}
Kasus \cmd{n=257} butuh 2 block; thread sisa pada block kedua ditangani guard batas dengan benar.

Modifikasi (\texttt{src/vector\_add\_alpha.cu}): kernel diubah menjadi \cmd{c[i] = 2*a[i] + b[i]}; untuk \cmd{n=1000003} hasilnya \cmd{errors=0}.

\section{Exercise 4: Waktu kernel vs waktu termasuk transfer}
Berkas: \texttt{src/vector\_benchmark.cu}. Program mengukur tiga hal dengan CUDA event: waktu kernel saja (\cmd{kernel}, hanya launch + eksekusi), waktu operasi penuh (\cmd{operation}: alokasi, dua H2D copy, kernel, satu D2H copy), dan waktu CPU sebagai pembanding. Dilaporkan dua speedup: \emph{resident-data speedup} = CPU/kernel (data sudah di device) dan \emph{operation speedup} = CPU/operation (termasuk transfer). Setiap ukuran dijalankan 5 kali; Tabel~\ref{tab:ex4} memakai median.

\begin{table}[H]\centering
\caption{Exercise 4 --- median 5 run per ukuran (Tesla T4).}
\label{tab:ex4}
\begin{tabular}{rrrrrr}
\toprule
\cmd{n} & CPU (ms) & kernel (ms) & operasi (ms) & speedup residu & speedup operasi\\
\midrule
1003 & 0.0004 & 0.0082 & 0.0443 & 0.05 & 0.01\\
100003 & 0.0344 & 0.0162 & 0.4099 & 2.13 & 0.08\\
1000003 & 0.5801 & 0.0525 & 3.0029 & 11.04 & 0.19\\
10000003 & 8.8437 & 0.4576 & 25.9355 & 19.33 & 0.34\\
\bottomrule
\end{tabular}
\end{table}

\plotfig{0.85\linewidth}{figures/vector_speedup.png}{Speedup GPU vs CPU terhadap ukuran vektor (median 5 run). Speedup kernel-only tumbuh hingga $\approx$19.3$\times$ pada $n=10^7$, tetapi speedup termasuk transfer tidak pernah melewati 1.}

Dua batas waktu (\emph{timing boundaries}) didefinisikan eksplisit: \cmd{kernel} hanya melingkupi launch kernel (diukur dengan CUDA event, tanpa transfer), sedangkan \cmd{operation} mencakup alokasi device, copy H2D/D2H, dan kernel. Hasilnya menegaskan pelajaran utama exercise ini: kernel saja 19$\times$ lebih cepat dari CPU pada $n=10^7$, tetapi sekali transfer PCIe ikut dihitung, operasi penuh justru $\approx$3$\times$ \emph{lebih lambat} dari CPU (speedup 0.34). GPU menang besar hanya bila data residen di device dan dipakai ulang.

\section{Exercise 5: Kernel grid-stride, tuning launch, bandwidth efektif}
Berkas: \texttt{src/vector\_gridstride.cu}. Kernel memakai pola \emph{grid-stride loop}: tiap thread memproses indeks \cmd{tid}, \cmd{tid+stride}, \cmd{tid+2*stride}, \ldots\ dengan \cmd{stride} sama dengan jumlah thread total, sehingga konfigurasi block/thread bebas dipilih tanpa mengubah kebenaran --- kernel tetap mencakup seluruh $n$ lewat stride. Bandwidth efektif dihitung sebagai $3 \times n \times 4$ byte (baca \cmd{a}, \cmd{b}, tulis \cmd{c}) dibagi waktu kernel.

Sweep 9 konfigurasi (64/128/256 block $\times$ 64/128/256 thread, 3 run tiap konfigurasi, median pada Tabel~\ref{tab:ex5} dan Gambar~\ref{fig:bw}):

\begin{table}[H]\centering
\caption{Exercise 5 --- median 3 run per konfigurasi, $n=1000003$.}
\label{tab:ex5}
\begin{tabular}{rrrr}
\toprule
block & thread & kernel (ms) & bandwidth efektif (GB/s)\\
\midrule
256 & 256 & 0.0530 & 226.45\\
128 & 256 & 0.0534 & 224.69\\
256 & 128 & 0.0547 & 219.43\\
64 & 256 & 0.0628 & 191.13\\
128 & 128 & 0.0638 & 187.97\\
256 & 64 & 0.0692 & 173.45\\
128 & 64 & 0.0986 & 121.71\\
64 & 128 & 0.1006 & 119.28\\
64 & 64 & 0.1677 & 71.55\\
\bottomrule
\end{tabular}
\end{table}

\plotfig{0.9\linewidth}{figures/gridstride_bw.png}{Bandwidth efektif terhadap konfigurasi launch (median 3 run, $n=1000003$). Konfigurasi 256 block $\times$ 256 thread tercepat (226 GB/s); 64 $\times$ 64 paling lambat (72 GB/s).}
\label{fig:bw}

Konfigurasi terbaik (256$\times$256, 226 GB/s) $\approx$3.2$\times$ lebih cepat dari yang terburuk (64$\times$64, 72 GB/s). Polanya jelas: thread block yang lebih besar (256) konsisten menang karena occupancy dan coalescing akses memori yang lebih baik; 65536 thread total (256$\times$256) memberi cukup paralelisme untuk menjenuhkan 40 SM pada T4.

\section{Exercise 6: Reduksi dengan shared memory}
Berkas: \texttt{src/reduce\_sum.cu}. Tiap block mereduksi 256 elemen memakai array \emph{shared memory} 256 elemen, pola halving dengan \cmd{\_\_syncthreads()} tiap tahap, lalu menjumlahkan hasil per-block di host. Untuk \cmd{n=1000003}: \cmd{expected=3000003 result=3000003 match=yes}. Uji ukuran tepi semuanya cocok:

\begin{center}
\begin{tabular}{rrrr}
\toprule
\cmd{n} & expected & result & match\\
\midrule
1 & 0 & 0 & yes\\
255 & 759 & 759 & yes\\
256 & 762 & 762 & yes\\
257 & 766 & 766 & yes\\
1000003 & 3000003 & 3000003 & yes\\
\bottomrule
\end{tabular}
\end{center}

\section{Exercise 7: Perkalian matriks naive vs tiled}
Berkas: \texttt{src/matrix\_mul.cu} (naive) dan \texttt{src/matrix\_mul\_tiled.cu} (tiled $16\times16$). Benchmark \texttt{src/matrix\_bench.cu} menguji keduanya; varian tiled dipilih via \cmd{-DUSE\_TILED}. Kebenaran dicek pada matriks non-persegi $127 \times 130 \times 129$: kedua varian \cmd{errors=0}, \cmd{max\_error=0.000e+00}.

Perbandingan waktu pada ukuran persegi (5 run per varian per ukuran, median pada Tabel~\ref{tab:ex7}, Gambar~\ref{fig:mm}):

\begin{table}[H]\centering
\caption{Exercise 7 --- median 5 run per varian per ukuran (Tesla T4).}
\label{tab:ex7}
\begin{tabular}{rrrr}
\toprule
varian & N & kernel (ms) & GFLOP/s\\
\midrule
naive & 128 & 0.0201 & 208.38\\
tiled & 128 & 0.0163 & 257.51\\
naive & 256 & 0.0996 & 336.95\\
tiled & 256 & 0.0655 & 512.00\\
naive & 512 & 0.8602 & 312.08\\
tiled & 512 & 0.5562 & 482.60\\
\bottomrule
\end{tabular}
\end{table}

\plotfig{0.95\linewidth}{figures/matmul_naive_tiled.png}{Waktu kernel (kiri, skala log) dan throughput (kanan) matmul naive vs tiled terhadap ukuran matriks (median 5 run).}
\label{fig:mm}

Varian tiled konsisten $\approx$1.5$\times$ lebih cepat (N=512: 0.556 ms vs 0.860 ms; 482.6 vs 312.1 GFLOP/s) karena ubin 16$\times$16 di shared memory memotong akses berulang ke memori global sebesar faktor 16. Seluruh 30 run benchmark \cmd{errors=0}.

\section{Exercise 8: Diagnosis error dengan Compute Sanitizer}
\cmd{compute-sanitizer --tool memcheck} tersedia pada CUDA 12.8 dan dijalankan pada \cmd{vector\_add}, \cmd{reduce\_sum}, dan \cmd{matrix\_mul\_tiled}: ketiganya \cmd{ERROR SUMMARY: 0 errors} --- tidak ada out-of-bounds atau akses ilegal pada program yang benar.

Dua bug terkontrol kemudian disuntikkan pada \emph{salinan} berkas (aslinya tidak disentuh):
\begin{enumerate}
\item \textbf{Tulis di luar batas} pada salinan \cmd{vector\_add} (indeks ditulis sejauh $n$ elemen melewati akhir buffer): sanitizer melaporkan \cmd{CUDA error: unspecified launch failure} dengan 97 error. Kernel yang menulis ke memori ilegal menggagalkan launch secara keseluruhan.
\item \textbf{Menghapus barrier pertama} pada salinan \cmd{reduce\_sum}: \cmd{--tool racecheck} melaporkan hazard baca/tulis di baris 16 dan hasil reduksi salah (\cmd{result=2644645}, \cmd{expected=3000003}, \cmd{match=no}). Tanpa barrier, thread membaca \cmd{sdata} yang belum selesai ditulis thread lain --- race condition klasik yang hanya muncul pada eksekusi paralel.
\end{enumerate}

\section{Mini-project: Stencil 2D lima titik}
Berkas: \texttt{src/stencil.cu}. Kernel menerapkan stencil lima titik (tengah + atas/bawah/kiri/kanan) pada grid 2D; guard menangani tepi domain. Dua mode waktu diukur: \emph{resident} (satu H2D, $k$ langkah kernel, satu D2H) dan \emph{per-step} (transfer bolak-balik tiap langkah). Hasil pada Tabel~\ref{tab:stencil}, seluruhnya \cmd{errors=0} (verifikasi terhadap CPU).

\begin{table}[H]\centering
\caption{Mini-project stencil --- median waktu (Tesla T4).}
\label{tab:stencil}
\begin{tabular}{lrrrr}
\toprule
kasus & CPU & kernel residu & operasi residu & per-step transfer\\
\midrule
$257 \times 259$, 10 langkah & 1.559 ms & 0.061 ms & 0.238 ms & 1.562 ms\\
$1024 \times 1024$, 20 langkah & 50.648 ms & 0.936 ms & 3.103 ms & 36.014 ms\\
\bottomrule
\end{tabular}
\end{table}

Pola Exercise 4 terulang dengan tajam: kernel residu 54$\times$ lebih cepat dari CPU pada $1024^2$, tetapi memindahkan data bolak-balik tiap langkah (36.0 ms) $\approx$12$\times$ lebih lambat dari strategi residu (3.1 ms). Untuk stencil iteratif, data \emph{harus} tinggal di device antar langkah.

\section{Refleksi: perlambatan, perbaikan, dan bug}
\textbf{Perlambatan --- transfer mendominasi.} Pada Exercise 4, speedup termasuk transfer tidak pernah melewati 1 (maksimum 0.34 pada $n=10^7$): operasi penuh justru $\approx$3$\times$ lebih lambat dari CPU meski kernelnya 19$\times$ lebih cepat. Stencil mengonfirmasi dari sisi lain: transfer per-langkah 12$\times$ lebih lambat dari eksekusi residu. Pelajarannya: bandwidth PCIe ($\approx$25 ms untuk 120 MB pada $n=10^7$) adalah biaya tetap yang hanya terbayar bila komputasi cukup besar atau data dipakai ulang.

\textbf{Perbaikan --- tiling dan tuning launch.} Tiled matmul mencapai 482.6 GFLOP/s vs 312.1 GFLOP/s naive pada N=512 (1.55$\times$) lewat pemakaian ulang data di shared memory. Pada grid-stride, konfigurasi 256$\times$256 (226 GB/s) mengalahkan 64$\times$64 (72 GB/s) sebesar 3.2$\times$ --- thread block besar memberi occupancy dan coalescing yang lebih baik pada 40 SM T4.

\textbf{Bug yang didiagnosis --- ekspansi variabel IPython.} Pada run pertama, lima cell loop shell (\cmd{!for n in ...; do ... \$n ...; done}) semuanya berjalan dengan $n=10000003$: Jupyter mengintersep \cmd{\$n} dan menggantinya dengan variabel \emph{Python} \cmd{n} (bernilai 10000003 dari cell ringkasan), bukan variabel loop shell. Gejalanya jelas pada output: lima baris identik \cmd{n=10000003} dan, pada cell benchmark matriks, peringatan \emph{integer overflow} diikuti \cmd{CUDA error: invalid configuration argument} karena ukuran matriks menjadi $10000003 \times 10000003$. Diagnosis: \cmd{!} menjalankan perintah lewat IPython yang melakukan ekspansi \cmd{\$}; perbaikannya mengganti \cmd{!} dengan magic \cmd{\%\%bash} sehingga \cmd{\$n} menjadi milik bash sepenuhnya. Setelah perbaikan, seluruh ukuran/konfigurasi berjalan benar (20 run Exercise 4, 27 run Exercise 5, 30 run Exercise 7).

\section{Reproduksi}
Seluruh pekerjaan direproduksi lewat notebook Kaggle pada repositori:
\begin{enumerate}
\item Unggah \texttt{cuda\_tutorial\_kaggle.ipynb} via \cmd{File} $\rightarrow$ \cmd{Import Notebook}.
\item Pada panel kanan, set \cmd{Accelerator} ke \cmd{GPU}.
\item \cmd{Run All}; unduh notebook yang sudah tereksekusi via \cmd{File} $\rightarrow$ \cmd{Download as .ipynb}.
\end{enumerate}
Struktur repositori: \texttt{src/} (13 berkas \cmd{.cu}/\cmd{.cuh} untuk semua exercise dan mini-project), \texttt{evidence/} (output notebook + CSV data mentah 77 pengukuran), \texttt{latex\_report/} (laporan ini + gambar), dan notebook Kaggle asli. Flag kompilasi seragam: \cmd{-std=c++17 -O3 -lineinfo -arch=native}. Jika \cmd{-arch=native} gagal pada GPU lain, ganti dengan \cmd{-arch=sm\_70} (P100) atau \cmd{-arch=sm\_75} (T4).

\section{Kesimpulan}
Kedelapan exercise dan mini-project stencil diselesaikan pada Tesla T4 (CC 7.5, CUDA 12.8) dengan seluruh uji kebenaran lolos, termasuk ukuran tepi. Pengukuran menunjukkan tiga pelajaran utama: (1) kernel GPU mencapai speedup 19$\times$ atas CPU untuk data residu, tetapi transfer PCIe membuat operasi end-to-end lebih lambat dari CPU bila data hanya dipakai sekali; (2) tiling shared memory memberi 1.55$\times$ pada matmul dan tuning launch memberi hingga 3.2$\times$ pada bandwidth; (3) Compute Sanitizer bersih pada program benar dan menangkap kedua bug terkontrol (out-of-bounds write dan race condition reduksi). Satu bug alur kerja nyata (ekspansi \cmd{\$n} oleh IPython) didiagnosis dan diperbaiki dengan \cmd{\%\%bash}; seluruh angka pada laporan ini berasal dari run setelah perbaikan.

\end{document}
